\begin{titlepage}
  \centering
  \vspace*{1.8cm}
  {\LARGE\bfseries Единое Естествознание (Книга~1):\\ Дискретная квантовая механика — Теоретические основания\par}
  \vspace{0.4cm}
  {\Large\bfseries От отсутствия внешней памяти у физической природы и проективного исчисления\\
  к границе применимости непрерывной квантовой механики,\\
  атомным и ядерным оболочкам, планетогенезу, молекулярной биофизике\\
  и теоретико-модельному пределу в Lean~4\par}
  \vspace{0.5cm}
  {\normalsize\itshape Unified Natural Science (Book~1: Theoretical Foundations of Discrete Quantum Mechanics in Lean~4)\par}
  \vspace{1.0cm}
  {\large Академическая теоретическая монография\par}
  \vspace{0.35cm}
  {\normalsize Все теоремы формально верифицированы в ядре Lean~4 (\texttt{lean/MtCarrier/}, \texttt{0 sorryAx}, \texttt{0 Classical.choice}).\\
  Лабораторные протоколы чистой комнаты (\texttt{M1}--\texttt{M4}, \texttt{F1}--\texttt{F6}) см.\ в \textbf{Книге~2} (\texttt{dqm-laboratory-manual.pdf});\\
  теоретическую электротехнику и кремниевый стек см.\ в \textbf{Книге~3} (\texttt{dqm-engineering-handbook.pdf});\\
  популярный путеводитель (\emph{Powered by Dr.\,STONE}) см.\ в \textbf{Книге~4} (\texttt{natural-sciences-course.pdf}).\par}
  \vfill
  {\large\bfseries Автор: \bookauthorfullname\par}
  \vspace{0.2cm}
  {\footnotesize\itshape Формальная верификация в Lean~4 и аналитическая разработка выполнены при участии ИИ-партнёра \textbf{Луч} (\textbf{Luch}, AI-Guiders).\par}
  \vspace{0.3cm}
  {\large \today\par}
\end{titlepage}
\clearpage

\chapter*{Условные обозначения}
\markboth{Условные обозначения}{Условные обозначения}
\addcontentsline{toc}{chapter}{Условные обозначения}

{\small
\renewcommand{\arraystretch}{1.30}
\setlength{\tabcolsep}{5pt}
\begin{longtable}{@{}p{0.21\textwidth}p{0.58\textwidth}p{0.15\textwidth}@{}}
\toprule
\textbf{Обозначение} & \textbf{Определение и физический смысл} & \textbf{Значение} \\
\midrule
\endfirsthead
\toprule
\textbf{Обозначение} & \textbf{Определение и физический смысл} & \textbf{Значение} \\
\midrule
\endhead
\bottomrule
\endfoot
\multicolumn{3}{@{}l}{\textit{I. Двусортная система носителей и фундаментальные константы при $d = 3$}} \\
\midrule
$\mathcal{M},\ \mathcal{T}$ & Дискретный носитель фазовых колец и непрерывный континуальный носитель & Два сорта \\
$\pi_{M\to T}$ и $\iota_{T\to M}$ & Проективный подъём Вейерштрасса и линейная ретракция носителей & Биекция \\
$d$ & Размерность пространства из уравнения самозамкнутости $1 + 2^d = d^2$ & $d = 3$ \\
$\mathrm{mod\_bits}(3)$ & Битовая разрядность однооборотного фазового регистра ячейки $d^2$ & $9\text{ бит}$ \\
$\phaseRingCount(3)$ & Модуль однооборотного фазового кольца ячейки $2^{\mathrm{mod\_bits}(3)}$ & $512 = 2^9$ \\
$\mu_{\mathrm{vac}}(3)$ & Неделимый вакуумный квант амплитуды (единица носителя $d^2 - 2^d$) & $1$ \\
$Q_\square(3)$ & Четверть-период фазового кольца и радиус алмаза $\phaseRingCount(3)/4$ & $128 = 2^7$ \\
$\Phi_{\mathrm{Pauli}}(3)$ & Полупериодный потолок антиподального запирания Паули $\phaseRingCount(3)/2$ & $256 = 2^8$ \\
$K_P(3)$ & Шаг биссектрисы стационарной фазы Френеля и мантисса $Q_\square(3)/2$ & $64 = 2^6$ \\
$\Delta\varphi_{\mathrm{seam}}(3)$ & Угловой шов одной грани кубооктаэдра $\lfloor \Phi_{\mathrm{Pauli}}(3) / 12 \rfloor$ & $21 = 3 \times 7$ \\
$\Delta\varphi_{\mathrm{disc}}(3)$ & Порог квантования (мёртвая зона Гейзенберга $(512+21)/13$) & $41$ \\
$|\mathcal{E}_{\mathrm{Heis}}(3)|$ & Число подпороговых состояний коридора Гейзенберга $2\Delta\varphi_{\mathrm{disc}}(3) - 1$ & $81 = 3^4$ \\
$B_1(3),\ B_2(3)$ & Однотактовый ($4\times 9$) и двухтактовый ($2\times 36$) битовые бюджеты ячейки & $36\text{ и }72\text{ бита}$ \\
$\mathrm{IV}(3)$ & Затравочное число Каталана $\Phi_{\mathrm{Pauli}}(3) + |{\triangle}(3)| - \mu_{\mathrm{vac}}(3)$ & $263$ \\
\midrule
\multicolumn{3}{@{}l}{\textit{II. Геометрия кубооктаэдра $O_2(3)$ и симплициального комплекса $\mathcal{C}_{13}$}} \\
\midrule
$O_2(3)$ & Оболочка кубооктаэдра из $12$ векторов $\mathbf{e}\in\{-1,0,+1\}^3$ с $e_x^2+e_y^2+e_z^2=2$ & $|O_2(3)| = 12$ \\
$N_\varphi(3)$ & Число узлов причинного ядра $B_1(3)$ ($1$ центр плюс $12$ соседей) & $13$ \\
$|\square(3)|,\ |{\triangle}(3)|$ & Числа квадратных ($2d$) и треугольных ($2^d$) граней кубооктаэдра & $6\text{ и }8$ \\
$d_{\mathrm{in}}(\square),\ d_{\mathrm{out}}(E,V)$ & Индекс квадратной грани и внешние индексы ребра и вершины в $O_2(3)$ & $4,\ 5,\ 7$ \\
$S_k(3)$ & Число узлов в $k$-й координационной сфере кубооктаэдра $10k^2 + 2$ & $12, 42, 92, 162$ \\
$V_{O_2}(3, N)$ & Число узлов в $N$-слойном шаре кубооктаэдра $1 + \sum_{k=1}^N (10k^2+2)$ & $13, 55, 147$ \\
$\mathcal{C}_{13}$ & Кликовый комплекс ядра $B_1(3)$ с числами клик $(13, 36, 32, 8)$ & $\chi(\mathcal{C}_{13}) = 1$ \\
$c_0$ и $\kappa_{\mathrm{geom}}$ & Скорость связи вдоль ребра $O_2(3)$ и геометрический фактор проекции & $\sqrt{2}\,c\text{ и }1/\sqrt{2}$ \\
\midrule
\multicolumn{3}{@{}l}{\textit{III. Дискретная экспонента Эйлера, оператор эволюции поля $\mathcal{G}_3$ и инварианты}} \\
\midrule
$\mathcal{C}_3(\varphi),\ \mathcal{S}_3(\varphi)$ & Чётная и нечётная проекции экспоненты Эйлера на $|\mathcal{C}_3|+|\mathcal{S}_3|=128$ & $[-128, +128]$ \\
$\mathcal{W}_3(\varphi)$ & Разностный вронскиан $\mathcal{C}_3(\varphi)\Delta\mathcal{S}_3(\varphi) - \mathcal{S}_3(\varphi)\Delta\mathcal{C}_3(\varphi)$ & $128 = Q_\square(3)$ \\
$\mathcal{G}_3$ и $\Psi^{(t)}$ & Оператор одного такта эволюции поля Верле и состояние на такте $t$ & $\texttt{fieldStepOp 3}$ \\
$q_e$ и $r_e$ & Частное переноса ветви и главный остаток разности фаз $\delta_e = 512 q_e + r_e$ & $q_e \in \{-1,0,1\}$ \\
$w_f$ и $\mathcal{B}_f$ & Вихревое число обмотки плакета $f$ и поток Максвелла $\mathcal{B}_f = 512\,w_f$ & $w_f \in \mathbb{Z}$ \\
$\hat{n}_E(\Phi)$ & Число излучённых надпороговых квантов энергии $\lfloor |\Phi| / 41 \rfloor$ & $\{0, \dots, 6\}$ \\
$\widehat{M}(3)$ & Простой гауссов кулоновский модуль решётки $4^2 + 9^2$ & $97$ \\
$\alpha_{\mathrm{fs}}^{-1}(3)$ & Четырёхпетлевая обратная постоянная тонкой структуры в $\mathbb{Q}(\sqrt{2})$ & $137{,}03599919$ \\
$\varepsilon(3)$ и $D(\mathbf{x})$ & Вакуумный пол импеданса ворот и защищённый импеданс ячейки & $\varepsilon = 8\,192$ \\
$N_{\mathrm{HB}}(3)$ & Делитель водородной связи $K_P(3) - \mu_{\mathrm{vac}}(3)$ для трёх окон биосферы & $63 = 7 \times 9$ \\
\end{longtable}
}
\clearpage
\tableofcontents
\clearpage
